\documentclass[10pt,a4paper]{amsart}
\usepackage[margin=19mm]{geometry}
\usepackage{amsmath,amssymb,mathtools}
\usepackage[scr=rsfs,cal=euler]{mathalfa}
\usepackage{xfrac}
\usepackage[unicode,hidelinks,bookmarks=false]{hyperref}
\newcommand{\ie}{i.e.\ }
\newcommand{\cf}{cf.\ }
\newcommand{\resp}{respectively\ }
\newcommand{\Subsection}{\subsection}
\providecommand{\nobreakdash}{\nobreak\hbox{-}}
\setlength{\emergencystretch}{2em}
\multlinegap0pt
\newtheorem{lemma}{Lemma}
\usepackage{cleveref}
\crefname{lemma}{Lemma}{Lemmas}
\newcommand{\set}[1]{\left\{#1\right\}}
\title{Infinitary positive existential normal forms for modules}
\author{Rishi Banerjee}
\date{Source: arXiv:2608.16692v3 (CC BY 4.0)}
\begin{document}
\maketitle

\noindent\emph{Source and licence.} Infinitary positive existential normal forms for modules. Version arXiv:2608.16692v3 (CC BY 4.0), distributed by arXiv under the Creative Commons Attribution 4.0 International licence, http://creativecommons.org/licenses/by/4.0/, as recorded in the arXiv licence metadata for this exact version and shown on the abstract page https://arxiv.org/abs/2608.16692v3. Full licence notice: \texttt{licenses/arxiv-2608.16692-CC-BY-4.0.txt}.\par\smallskip

\noindent\textit{Editorial scope.} Counted unit: the article's lemma lem:proj-atom (source0.tex lines 203--214). The article's complete original proof of it is delivered with it (source0.tex lines 216--302, one \texttt{proof} environment of 87 lines). No mathematics is written, repaired or re-derived: the statement and the proof are the article's own lines, in the article's own notation, and no retained line is substituted or re-flowed.  Retained uncounted as the source-local context the proof needs: lines 161--164.  Nothing was omitted from any retained span, so the package carries no omission record.  No retained line contains a citation, so the wrapper carries no bibliography and no bibliography entry.  No retained line contains a figure environment, an image-inclusion command or drawing code, so no image is delivered and the package declares no asset dependency.  Editorial formatting only, in addition: an AMS \texttt{amsart} preamble that restates the article's own declarations and macros used by the retained lines, namely the environments \texttt{lemma} on separate counters, together with the standard packages those lines need.  The retained environments print in this document as Lemma 1 in order of appearance, whereas the article prints the same items with the numbering of its own full document.  The complete native source of the article as distributed in the arXiv e-print for this exact version is bundled unchanged as \texttt{sources/207845/source0.tex}.
\begin{lemma}
\label{lem:test-set}
    Let $\kappa$ be an infinite cardinal, let $G$ be an abelian group, and let $(G_s)_{s \in S}$ be a family of subgroups of $G$ with $|S| \leq \kappa$. Then there is a test set $T$ for $(G, (G_s)_{s \in S})$ such that $|T| \leq 2^\kappa$.
\end{lemma}

\begin{lemma}
\label{lem:proj-atom}
    Let $\kappa$ be infinite and let 
    \[ Q = \bigcap_{i \in I} C_i \setminus \bigcup_{j \in J} D_j, \]
    where the $C_i$ and $D_j$ are cosets of subgroups $P_i, N_j \leq M^{\epsilon+ \zeta}$ and $|I|, |J| \leq \kappa$.
    
    If $Q \neq \emptyset$, then there is $a \in M^{\epsilon + \zeta}$, an index set $K$ with $|K| \leq 2^\kappa$, a family $(b_k)_{k \in K}$ where each $b_k \in M^{\epsilon + \zeta}$, and a family $(J_k)_{k \in K}$ where each $J_k \subseteq J$, such that 
    \[ 
        \pi(Q) = \left( \pi(a) + \pi(P) \right) \setminus \bigcup_{k \in K}  \left[ \pi(b_k) + \pi \left( P \cap \bigcap_{j \in J_k} N_j \right) \right],
    \]
    where $P = \bigcap_i P_i$.
\end{lemma}

\begin{proof}
    Let $C = \bigcap_{i \in I} C_i$, and note that $C$ is a coset of $P$. 
    Pick any point $a \in C$. 
    Then $\pi(C) = \pi(a) + \pi(P)$.
    
    Let $x \in \pi(C)$, and let $C_x, (D_j)_x \subseteq M^\zeta$ be the fibers over $x$:
    \begin{align*}
        C_x  = \set{ y \in M^\zeta \mid (x,y) \in C }, \qquad
        (D_j)_x  = \set{ y \in M^\zeta \mid (x,y) \in D_j }.
    \end{align*}
    Then 
    \[
        x \not \in \pi(Q) \iff C_x \subseteq \bigcup_{j \in J} (D_j)_x.
    \]
    Let $G, G_j$ be the vertical subgroups of $P, N_j$:
    \begin{align*}
        G = \set{ g \in M^\zeta \mid (0,g) \in P }, \qquad
        G_j = \set{ g \in M^\zeta \mid (0,g) \in N_j }.
    \end{align*}
    For any $y \in C_x$, $C_x - y = G$ and $(D_j)_x - y$ is empty or a coset of $G_j$.
    Let 
    \[
        B_j(x,y) = \set{ g \in G \mid y + g \in (D_j)_x }.
    \]
    For each $j$, $B_j(x,y)$ is either empty or a coset of $G \cap G_j$, and
    \[
        C_x \subseteq \bigcup_{j \in J} (D_j)_x \iff G \subseteq \bigcup_{j \in J} B_j(x,y).
    \]
    Applying \cref{lem:test-set} with $S = J$ to $(G, (G \cap G_j)_{j \in J})$, we obtain a test set $T \subseteq G$ with $|T| \leq 2^\kappa$ and 
    \[ 
        G \subseteq \bigcup_{j \in J} B_j(x,y) \iff T \subseteq \bigcup_{j \in J} B_j(x,y). \]
    Note that 
    \[ 
        t \in B_j(x,y) \iff (x,y) \in D_j - (0,t).
    \]
    So we have
    \begin{align*}
        x \not \in \pi(Q) 
        & \iff \exists y \in C_x \ \left[ T \subseteq \bigcup_{j \in J} B_j(x,y) \right] \\
        & \iff \exists y \in C_x \ \exists f: T \to J \ \forall t \in T \ \left[ (x,y) \in D_{f(t)} - (0,t) \right] \\
        & \iff \exists y \in C_x \ \exists f: T \to J \ \left[ (x,y) \in \bigcap_{t \in T} (D_{f(t)} - (0,t)) \right] \\
        & \iff x \in \bigcup_{f: T \to J} \pi \! \left( C \cap \bigcap_{t \in T} (D_{f(t)} - (0,t)) \right).
    \end{align*}
    For any $f: T \to J$, let
    \[ E_f = C \cap \bigcap_{t \in T} (D_{f(t)} - (0,t)), \]
    so the above becomes
    \[
        x \not \in \pi(Q) \iff x \in \bigcup_{f: T \to J} \pi(E_f).
    \]
    Each $E_f$ is either empty or a coset of $P \cap \bigcap_{t \in T} N_{f(t)} = P \cap \bigcap_{j \in f(T)} N_j$.
    For each nonempty $E_f$, pick an element $b_f \in E_f$, so that 
    \[  
        E_f = b_f + \left( P \cap \bigcap_{j \in f(T)} N_j \right).
    \]
    This gives us the following normal form for $\pi(Q)$:
    \begin{align*}
        \pi(Q) = \pi(C) \setminus \bigcup_{\substack{f: T \to J \\ E_f \neq \emptyset}} \left[ \pi(b_f) + \pi\! \left( P \cap \bigcap_{j \in f(T)} N_j \right) \right].
    \end{align*}

    There is one final bit of cardinality book-keeping to attend to. 
    We promised at most $2^\kappa$ many points $b_f$ and subsets $J_f$, but there may be more than $2^\kappa$ functions $T \to J$. However, we claim that there are only $2^\kappa$ many distinct nonempty $E_f$'s. 
    To see this, first rewrite $E_f$ as
    \[ 
        E_f = C \cap \bigcap_{t \in T} (D_{f(t)} - (0,t)) = C \cap \bigcap_{j \in f(T)} \bigcap_{t \in f^{-1}(j)} (D_j - (0,t)).
    \]
    Now if $f(t) = f(t') = j$, then $D_j - (0,t)$ and $D_j - (0,t')$ are both cosets of $N_j$.
    And if $E_f \neq \emptyset$, then these two cosets intersect, so they must be equal.
    So it suffices to pick for each $j \in f(T)$ a single $t_j \in f^{-1}(j) \subseteq T$. Define $\tau: J \to T \cup \set{\uparrow}$ by 
    \[
        \tau(j) = \begin{cases}
            t_j, & j \in f(T) \\
            \uparrow, & j \not \in f(T)
        \end{cases}.
    \]
    Then
    \[ 
        E_f = C \cap \bigcap_{j \in \tau^{-1}(T)} (D_j - (0,\tau(j))). 
    \]
    Thus every nonempty $E_f$ is determined by some map $\tau: J \to T \sqcup \set{\uparrow}$, and the number of such maps is at most 
    \[ (|T| + 1)^{|J|} \leq (2^\kappa)^\kappa = 2^\kappa. \]
    
    So, we may index the distinct nonempty sets $E_f$ as $\set{E_k \mid k \in K}$, where $|K| \leq 2^\kappa$. For each $k \in K$ choose a representative map $f_k: T \to J$ such that $E_k = E_{f_k}$. Choose $b_k \in E_{f_k}$ and put $J_k = f_k(T) \subseteq J$. Then 
    \[ 
        E_k = E_{f_k} = b_k + \left( P \cap \bigcap_{j \in J_k} N_j \right),
    \]
    and we have the desired normal form.
\end{proof}

\end{document}
